\chapter{ARIMA e Raiz Unitária}
\label{cap:arima}
% ── índice ──────────────────────────────────────────────────────
\index{ARIMA|textbf}
\index{raiz unitária|textbf}
\index{Dickey-Fuller, teste de|textbf}
\index{ADF|see{Dickey-Fuller, teste de}}
\index{adf@\texttt{adf()}|textbf}
\index{integração de ordem}
\index{KPSS, teste de}
\index{kpss@\texttt{kpss()}}

% ─────────────────────────────────────────────────────────────────────────────
\section{Processos integrados}

Um processo $\{y_t\}$ é \textbf{integrado de ordem $d$} — denotado $I(d)$
— quando precisa ser diferenciado $d$ vezes para se tornar estacionário.
O caso mais comum em economia é $I(1)$: a primeira diferença é estacionária,
mas o nível não.

O caso mais simples de processo $I(1)$ é o \textbf{passeio aleatório}
(\emph{random walk}):
\[
  y_t = y_{t-1} + \varepsilon_t, \qquad \varepsilon_t \sim \mathrm{WN}(0, \sigma^2)
\]

A variância cresce com $t$ ($\mathrm{Var}[y_t] = t\sigma^2$), de modo que
o processo não é estacionário. A primeira diferença $\Delta y_t = y_t -
y_{t-1} = \varepsilon_t$ é ruído branco — portanto $I(0)$.

O passeio aleatório \emph{com deriva} adiciona uma constante $\mu$:
\[
  y_t = \mu + y_{t-1} + \varepsilon_t
\]
com tendência determinística $E[y_t] = y_0 + \mu t$ — modelo básico para
preços de ativos, câmbio e muitas séries macroeconômicas.

% ─────────────────────────────────────────────────────────────────────────────
\section{Teste de Dickey-Fuller Aumentado}

O teste ADF (\emph{Augmented Dickey-Fuller}) testa formalmente a presença de
raiz unitária. O modelo auxiliar estimado é:
\[
  \Delta y_t = \alpha + \delta y_{t-1}
               + \sum_{j=1}^{k} \gamma_j \Delta y_{t-j} + \varepsilon_t
\]

com hipóteses:
\[
  H_0: \delta = 0 \quad \text{(raiz unitária)} \qquad
  H_1: \delta < 0 \quad \text{(estacionário)}
\]

A estatística $t$ de $\hat\delta$ não segue a distribuição $t$ padrão sob
$H_0$ — usa-se a distribuição de Dickey-Fuller com valores críticos
tabelados.

\begin{aviso}
  Os valores críticos do ADF são \textbf{negativos}: rejeita-se $H_0$
  quando a estatística é suficientemente \emph{negativa}. Um valor positivo
  (ou próximo de zero) não rejeita $H_0$ — evidência de raiz unitária.
\end{aviso}

% ─────────────────────────────────────────────────────────────────────────────
\section{O modelo ARIMA}

O ARIMA($p$,$d$,$q$) generaliza o ARMA para séries não estacionárias,
aplicando $d$ diferenciações antes da estimação:

\[
  \underbrace{\phi(L)(1-L)^d y_t}_{\text{AR em diferenças}}
  = c + \underbrace{\theta(L)\varepsilon_t}_{\text{MA}}
\]

onde $L$ é o operador de defasagem ($Ly_t = y_{t-1}$) e $(1-L)^d$ é o
operador de diferença de ordem $d$.

\begin{center}
\begin{tabular}{lll}
\toprule
\textbf{Modelo} & \textbf{Notação} & \textbf{Processo} \\
\midrule
Passeio aleatório & ARIMA(0,1,0) & $\Delta y_t = c + \varepsilon_t$ \\
AR(1) em diferenças & ARIMA(1,1,0) & $\Delta y_t = c + \phi\,\Delta y_{t-1} + \varepsilon_t$ \\
MA(1) em diferenças & ARIMA(0,1,1) & $\Delta y_t = c + \varepsilon_t + \theta\,\varepsilon_{t-1}$ \\
\bottomrule
\end{tabular}
\end{center}

% ─────────────────────────────────────────────────────────────────────────────
\section{Verificação por DGP}

\subsection*{Random walk: ADF não rejeita, AR(1): ADF rejeita}

O DGP contrasta os dois casos: passeio aleatório (I(1)) gerado por
\lstinline{cumsum(e)} versus AR(1) estacionário ($\phi=0{,}7$) gerado por
laço.

\begin{lstlisting}[caption={ADF em passeio aleatório vs.\ AR(1)}]
seed(42)
input df
id
1
end
for i in 1..=9 { df = append(df, df) }
df |> filter(_n <= 500)
generate df e  = rnormal()
generate df rw = cumsum(e)
generate df y  = 0.0
generate df t  = _n
let n = nrow(df)
let i = 2
while i <= n {
    let ly = mean(df, y, if = t == i - 1)
    let et = mean(df, e, if = t == i)
    replace df y = 0.7 * ly + et if t == i
    i = i + 1
}
adf(df, rw)
adf(df, y)
\end{lstlisting}

\begin{lstlisting}[language={}, basicstyle=\ttfamily\small,
  frame=none, numbers=none, backgroundcolor=\color{codbg},
  xleftmargin=0pt, xrightmargin=0pt, breaklines=false]
=============== Augmented Dickey-Fuller Test ===============
  Variable: rw   Lags used: 7
  H₀: series has a unit root (non-stationary)
  Test statistic:     -0.0980
  Critical values:  1%=-3.430  5%=-2.860  10%=-2.570
  Conclusion: Não rejeita H₀ — raiz unitária presente

=============== Augmented Dickey-Fuller Test ===============
  Variable: y   Lags used: 7
  H₀: series has a unit root (non-stationary)
  Test statistic:     -6.4616
  Critical values:  1%=-3.430  5%=-2.860  10%=-2.570
  Conclusion: REJEITA H₀ — estacionária
\end{lstlisting}

O passeio aleatório tem estatística $-0{,}10$ — acima dos valores críticos
negativos, não rejeita $H_0$. O AR(1) com $\phi=0{,}7$
tem estatística $-6{,}46$ — muito abaixo do valor crítico de 1\% ($-3{,}43$),
rejeita $H_0$ com alta confiança.

\subsection*{Estimação ARIMA(0,1,0) e seleção de ordem}

Com a raiz unitária identificada, estima-se com $d=1$. Três ordens
concorrentes para o passeio aleatório verdadeiro:

\begin{lstlisting}[caption={Seleção de ordem para série I(1)}]
let m010 = arima(df, rw, p=0, d=1, q=0)
let m110 = arima(df, rw, p=1, d=1, q=0)
let m011 = arima(df, rw, p=0, d=1, q=1)
display m010
display m110
display m011
\end{lstlisting}

\begin{lstlisting}[language={}, basicstyle=\ttfamily\small,
  frame=none, numbers=none, backgroundcolor=\color{codbg},
  xleftmargin=0pt, xrightmargin=0pt, breaklines=false]
-- ARIMA(0,1,0): AIC=156.77  BIC=160.97
   intercept  0.4863  std=0.0128  z=38.14  p=0.000

-- ARIMA(1,1,0): AIC=158.70  BIC=167.10
   intercept  0.4805  std=0.0253  z=18.96  p=0.000
   ar.L1      0.0120  std=0.0450  z=0.266  p=0.791  [insig.]

-- ARIMA(0,1,1): AIC=157.80  BIC=166.19
   intercept  0.4855  std=0.0128  z=38.07  p=0.000
   ma.L1      0.0058  std=0.0453  z=0.128  p=0.898  [insig.]
\end{lstlisting}

O modelo verdadeiro ARIMA(0,1,0) tem o menor AIC e BIC. Os termos AR e MA
adicionais são insignificantes ($p > 0{,}7$) — o BIC, que penaliza mais
fortemente a complexidade, confirma a parcimônia do modelo simples.

% ─────────────────────────────────────────────────────────────────────────────
\section{Fluxo de trabalho}

O protocolo padrão para séries temporais univariadas:

\begin{enumerate}
  \item \textbf{Inspecionar} o nível da série — tendência visível sugere
        não-estacionariedade.
  \item \textbf{Testar} com ADF — não rejeitar $H_0$ indica $d \geq 1$.
  \item \textbf{Diferenciar} uma vez e re-testar com ADF na série
        diferenciada — confirma $d=1$.
  \item \textbf{Selecionar} $(p,q)$ pela grade AIC/BIC sobre a série
        diferenciada.
  \item \textbf{Validar} resíduos: devem ser ruído branco.
\end{enumerate}

% ─────────────────────────────────────────────────────────────────────────────
\section{Sintaxe}

\begin{lstlisting}
adf(df, y)
adf(df, y, lags=4)
arima(df, y, p=0, d=1, q=0)
arima(df, y, p=1, d=1, q=1)
\end{lstlisting}

\begin{lstlisting}[caption={Fluxo típico com dados reais}]
load "cambio.csv" as df

adf(df, taxa)

// se não rejeitar H0: estimar com d=1
let m = arima(df, taxa, p=1, d=1, q=0)
display m
let yhat = predict(m, "xb")
\end{lstlisting}
